\chapter{Introduction}
\label{chap:introduction}

\section{Purpose}

\texttt{sitas} is an experimental Rust runtime and service model inspired by
Seastar's shard-per-core architecture. It explores shard ownership, explicit
cross-shard transfer, custom future execution, and operating-system I/O without
depending on a third-party async runtime.

The project started as a standard-library-only sharded key-value store. It is
now a four-crate workspace with a \texttt{no\_std + alloc} runtime contract,
host and CharlotteOS backends, a custom executor, Unix readiness, TCP and UDP
helpers, shard-local values, observability, CPU and memory placement, and Linux
\texttt{io\_uring} support.

\section{Prerequisites}

This manual assumes familiarity with:

\begin{itemize}
    \item Rust ownership, borrowing, and lifetimes;
    \item \texttt{Future}, \texttt{Poll}, \texttt{Waker}, and \texttt{Pin};
    \item basic concurrency: threads, channels, mutexes;
    \item Unix file descriptors and non-blocking I/O (for the I/O chapters).
\end{itemize}

\section{How to read this project}

Application code sees typed, owned boundaries. A caller asks a shard to do
work. The owning shard mutates its local state and returns an owned reply. The
implementation avoids hiding service state behind a global mutex.

The runtime keeps that model honest:

\begin{enumerate}
    \item futures are owned by executor tasks;
    \item tasks are assigned to one executor shard;
    \item wakers only requeue tasks on that executor;
    \item cross-shard work must be submitted explicitly;
    \item shard-local values can be borrowed only inside synchronous closures;
    \item observability returns copies, not borrowed internals.
\end{enumerate}

The following chapters move from ownership rules to task execution, sharding,
and I/O. Later chapters provide examples, source-reading paths, labs, and
validation commands.

\begin{notabene}[Learning goal]
You understand the core design when you can explain where a future is polled,
who owns the state it mutates, what wakes it, and what happens when it is
dropped.
\end{notabene}

\section{Why the project has this shape}

The runtime keeps task allocation, reactor registration, cross-thread wakeups,
cancellation, and kernel ownership visible. This makes the implementation less
polished than a general-purpose runtime, but easier to inspect.

\texttt{sitas} therefore accepts some explicitness that a general-purpose
runtime might smooth over:

\begin{description}
    \item[Explicit shards] make ownership visible. If work runs on shard 2,
        shard 2 owns the mutation.
    \item[Explicit submitters] make cross-shard movement visible. A future does
        not silently migrate just because another shard is less busy.
    \item[Explicit snapshots] make observation visible. Monitoring is a read of
        owned runtime facts, not a borrowed view into live internals.
    \item[Explicit OS boundaries] make platform behavior visible. A readiness
        wait, an \texttt{io\_uring} completion, and a CPU affinity request are
        different contracts with the operating system.
\end{description}

The project therefore prefers explicit boundary-crossing APIs.

\section{Scope of this manual}

This manual is the printable, long-form companion to the Markdown architecture
notes.

It focuses on:

\begin{itemize}
    \item architectural invariants;
    \item executor and scheduling behavior;
    \item shard ownership and cross-shard submission;
    \item readiness, \texttt{epoll}, and \texttt{io\_uring} integration;
    \item example programs that exercise the runtime;
    \item build and validation operations.
\end{itemize}

\begin{notabene}
The API is experimental and may change as the runtime contracts become clearer.
\end{notabene}

\section{What this project is not}

\texttt{sitas} is neither a production replacement for mature Rust runtimes
nor a line-by-line Seastar port. Its mechanisms remain small enough to inspect:

\begin{itemize}
    \item no Tokio, async-std, Glommio, Monoio, or actor framework is used as
          the core runtime;
    \item blocking standard-library service APIs remain distinct from
          awaitable executor APIs;
    \item unsafe and FFI code is kept near the operating-system boundary;
    \item mechanisms are introduced only after their semantics can be tested.
\end{itemize}

The resulting code demonstrates timers, wakers, readiness, \texttt{io\_uring},
shard-local values, scheduling groups, and placement in ordinary Rust modules.

\section{Further reading}

\begin{description}
    \item[Seastar] The C++ shard-per-core framework that inspired the
        ownership model.\\
        \url{https://seastar.io}
    \item[ScyllaDB shard-per-core architecture] Practical discussion of
        shared-nothing database design on top of Seastar.\\
        \href{https://www.scylladb.com/product/technology/shard-per-core-architecture/}{ScyllaDB technology overview}
    \item[Glommio] A Rust thread-per-core runtime using \texttt{io\_uring},
        useful for comparing design choices.\\
        \url{https://github.com/DataDog/glommio}
\end{description}
